\documentclass{jsarticle}
\usepackage{amsmath}
\usepackage{amssymb}
\begin{document}
\title{UFOでの準同型写像の回転体でのブラックホール \\
重力場エネルギーに対する電磁気力}
\author{Masaaki Yamaguchi}
\date{}
\maketitle

  $$ ||ds^2|| = 8 \pi G\left({{p \over c^3} + {V \over S}}\right) $$
  宇宙での位置座標と運動エネルギー、準同型写像で、宇宙のブラックホール
  の熱エントロピー値
  $$ \bigoplus {(i \hbar^{\nabla})}^{\oplus L} $$
  $$ = H \Psi = \nu \lambda $$
  UFOの位置とエネルギー関係式
  $$ {\pi{{(\int \Gamma(\gamma)^{'})^2}dx_m} \over \nabla M} $$
  種数に対する基本群
  $$ \pi(\chi) = {d \over dM}L$$
  その回転体での準同型写像でのブラックホールの重力場エネルギー
  $$ = {\pi{(\int {{\Gamma(\gamma)^{'}}}^{\nabla})}}^{2M}dx_m $$
  その電磁気力
  $$ = \int e^{-x \log x}\mathrm{div}(\mathrm{rot} E)dx $$
    
  鶏が先か卵が先かの違い。
  空間に穴が開くから、ブラックホールがあると言われている。
  リッチ・フローによる空間分類が、空間の種数によって、
  その中で生成されうる確率で、どの物体が存在、生成されるか決まる。
  これらより、その時空に関わる時間の性質が決まる。
  時間の光速度不変に対する遅れ、早まり(タキオン無いと言われている）、
  密度エネルギーの宇宙の周空間で説明できる。空間による進度の方向、
  どの物体が、どの種類の空間に属するのが適切かが決まる。

  量子コンピュータの重ね合わせが$\oplus L$で、
  本体が、$\bigoplus(i\hbar^{\nabla})$
  これが、ゼータ関数の使われ方であり、
  $\oplus L$が、Morse理論での種数のバンドル体であり、
  これが、目として、ブレインインターフェースに使われる。

  ソレノイドコイルの応用が、メビウス回路の反重力発生に使われる。
  電場・磁場を発生させる。
  キルヒホッフの法則で、抵抗値で電圧を上げたり、下げたりさせる。
  これも、種数の重ね合わせで実現できる。
  D-waveで、周期を止めて、計る。

  これらの理論は、空間概念の量子化で理論展開される。

\end{document}
